% !TEX root = ../../main.tex

\chapter{神经网络优化算法}

我们现在已经学会了用反向传播来计算梯度，并利用梯度下降来更新模型参数。一切看起来都很完美，数学推导无懈可击，似乎我们只需要把数据喂给模型，模型就能轻而易举地找到最优解。然而，真实的深度学习训练过程并没有这么简单。

在理想情况下，我们希望损失函数的地形如\cref{fig:convex_landscape}所示，呈现一个平滑的“U型碗”状。这样的函数属于\textbf{凸函数（Convex Function）}。对于凸函数，任意局部最优解都是全局最优解，因此梯度下降能够沿着损失下降的方向不断逼近最低点，并较为稳定地收敛到全局最优解。
\begin{figure}[H]
    \boxedimage[width=0.4\textwidth, trim=0 {0.2\height} {0.01\width} 0, clip]{figures/part_01/p1_ch03/convex_landscape.pdf}
    \caption{凸函数的损失曲面。曲面呈平滑的“U型碗”形态，任何局部极小值均为全局最小值。}
    \label{fig:convex_landscape}
\end{figure}

然而，真实的深度学习模型通常包含数万乃至上亿个参数，其对应的损失函数往往是复杂的\textbf{非凸函数（Non-convex Function）}，如\cref{fig:nonconvex_landscape}所示。与理想的“U型碗”不同，非凸损失地形往往更加复杂，可能包含局部极小值、鞍点以及高曲率区域。在这样的地形中，梯度下降不再能够轻松地找到全局最优解，参数更新可能在陡峭区域来回震荡，也可能在平坦区域前进缓慢，从而使优化过程变得十分困难。
\begin{figure}[H]
    \boxedimage[width=0.5\textwidth, trim={0.05\width} {0.18\height} {0.05\width} {0.12\height}, clip]{figures/part_01/p1_ch03/nonconvex_landscape.pdf}
    \caption{非凸函数的损失曲面。曲面拓扑结构复杂，充斥着局部极小值、鞍点与高曲率区域。}
    \label{fig:nonconvex_landscape}
\end{figure}

因此，仅靠最基础的梯度下降往往难以满足实际训练的需求。我们需要对这一参数更新公式$\theta\leftarrow\theta-\eta\frac{\partial\mathcal{L}}{\partial\theta}$进行进一步改进，使模型能够更加高效、稳定地优化参数。

\section{学习率}

我们首先回顾一下梯度下降的核心公式：
$$
\theta \leftarrow \theta - \eta \cdot \frac{\partial \mathcal{L}}{\partial \theta}
$$
在这个公式中，$\frac{\partial \mathcal{L}}{\partial \theta}$表示损失函数对当前参数的梯度，它指向损失函数增长最快的方向；而负梯度$-\frac{\partial \mathcal{L}}{\partial \theta}$就指向损失函数下降最快的方向，也是参数更新的方向。学习率$\eta$是步长系数，用来缩放梯度；参数每次实际更新的幅度由学习率与梯度大小共同决定。因此可以直观地理解为：负梯度决定“往哪个方向走”，学习率决定“沿这个方向走多远”。

直觉上，学习率设置得过大或者过小，都可能给模型训练带来问题，正所谓过犹不及。那么，学习率设置得过大或者过小，究竟会有什么后果呢？

\cref{fig:learning_rate_bad}左图展示了学习率过大的情况。由于每次参数更新的幅度过大，参数很容易一次跨过Loss较低的位置，直接进入另一次梯度（$\frac{\partial \mathcal{L}}{\partial \theta}$）更大的区域，从而导致下一次更新的幅度进一步增大。如此反复，参数很可能在最优解两侧不断来回越过，形成剧烈震荡（Oscillation），甚至逐渐发散，使Loss持续增大，最终因数值溢出而出现\texttt{NaN}（Not a Number），导致训练失败。

\begin{figure}[H]
    \boxedimage[width=0.9\textwidth, trim=0 0 0 0, clip]{figures/part_01/p1_ch03/learning_rate_bad.pdf}
    \caption{学习率过大与过小的对比。左图学习率过大，参数在损失曲线两侧震荡甚至发散，无法收敛；右图学习率过小，参数步长极小，向局部最小值的收敛极其缓慢。}
    \label{fig:learning_rate_bad}
\end{figure}

\cref{fig:learning_rate_bad}右图展示了学习率过小的情况。此时，每次参数更新的幅度都非常有限，参数虽然会沿着负梯度方向缓慢移动，但需要经过大量迭代才能明显降低Loss，导致训练收敛速度过慢。更重要的是，在复杂的损失函数中，参数可能会遇到局部极小值或鞍点（Saddle Points）。而由于局部极小值和鞍点附近的梯度往往接近于$0$，当学习率本身也很小时，$\eta$与$\frac{\partial\mathcal{L}}{\partial\theta}$共同作用，会使参数每次更新的幅度进一步减小。模型一旦进入这些区域，可能在整个训练周期内都无法离开，从而使训练速度显著下降，甚至出现近似停滞的现象。

\begin{orangeBox}{鞍点}
鞍点是损失函数中的一种特殊位置：它的梯度为零，但既不是局部极小值，也不是局部极大值。

\begin{figure}[H]
    \boxedimage[width=0.4\textwidth, trim=0 {0.08\height} 0 {0.1\height}, clip]{figures/part_01/p1_ch03/saddle_landscape.pdf}
    \caption{经典马鞍面$z = x^2 - y^2$。曲面沿$x$轴方向向上弯曲、沿$y$轴方向向下弯曲，形似马鞍。红点为鞍点：该处梯度为零，既非局部极小值也非局部极大值。}
    \label{fig:saddle_landscape}
\end{figure}

在鞍点附近，沿着某些方向移动，Loss会增大；而沿着另一些方向移动，Loss反而会减小。它既不像局部极小值：在局部极小值附近，向任意方向移动，Loss都会增大；也不像局部极大值：在局部极大值附近，向任意方向移动，Loss都会减小。如\cref{fig:saddle_landscape}所示，这种曲面的形状看起来很像马鞍，所以该位置的点被称为“鞍点”。
\end{orangeBox}

理想的训练过程应该如\cref{fig:learning_rate_good}所示。在训练初期，参数距离最优解还比较远，此时可以采用较大的学习率，让参数以较大的步长快速向最优解靠近；随着训练进行，参数逐渐接近最优解，此时如果仍然保持较大的步长，就容易在最优解附近来回震荡。因此，训练后期需要逐渐减小学习率，让参数以更小的步长进行调整，从而更加稳定地逼近最优解。

\begin{figure}[H]
    \boxedimage[width=0.45\textwidth, trim=0 0 0 0, clip]{figures/part_01/p1_ch03/learning_rate_good.pdf}
    \caption{合理设置学习率下的梯度下降轨迹。在远离全局最小值时采用较大学习率，参数得以快速下降并越过较浅的局部极小值（local minimum）；当接近全局最小值时逐渐减小学习率，避免在谷底附近震荡，最终平稳、精确地收敛至全局最小值（global minimum）。}
    \label{fig:learning_rate_good}
\end{figure}

既然训练的不同阶段对学习率有不同的需求，那么学习率就应该随着训练过程动态调整：训练初期保持较大，训练后期逐渐减小。

为了实现学习率随训练过程逐渐减小，一个直接的做法，就是每隔固定的训练轮次，将当前学习率乘以一个小于$1$的系数，使学习率在特定时刻快速下降。例如，每训练完$10$个Epoch，就将学习率缩减为原来的十分之一。这种按照固定间隔逐次降低学习率的策略，就是\textbf{阶跃衰减（Step Decay）}。其数学表达式为：
$$
\eta_t = \eta_0 \cdot \gamma^{\left\lfloor \frac{t}{s} \right\rfloor}
$$
其中，$\eta_0$表示初始学习率，$\gamma$表示衰减系数，通常取$0.1$，$t$表示当前的训练轮次（Epoch），$s$表示衰减步长，即每经过多少个Epoch进行一次衰减。例如，当$s=10$时，模型每训练$10$个Epoch，学习率就衰减一次。

可以看到，在这种策略下，学习率在一段时间内保持不变，到达设定的Epoch后突然下降一个台阶。训练初期，模型使用较大的学习率快速调整参数；随着训练进行，学习率逐步降低，使参数更新越来越精细。不过，这种方法也存在明显的不足：学习率的变化并不连续，而是呈阶梯状下降。每次学习率突然降低后，Loss的下降速度也可能随之发生明显变化，使整个训练过程显得不够平滑。

为了让学习率能够更加平滑地下降，我们希望它随着训练过程连续变化，而不是在某几个时刻突然降低。同时，学习率的变化速度也应该有所区别：训练初期下降得比较缓慢，以保留较大的学习率进行快速调整；训练中期下降速度加快，使学习率逐渐减小；到了训练后期再次放缓，让模型以更小的步长进行精细调整。这种变化方式与整个训练过程比较契合。那么，有没有一种函数能够满足这样的变化规律呢？如\cref{fig:cosine}所示，余弦函数$\cos(x)$在$[0,\pi]$区间内恰好具有这样的特点：两端变化较慢，中间变化较快，因此可以用它来构造平滑的学习率衰减过程。

\begin{figure}[H]
    \boxedimage[width=0.45\textwidth, trim=0 0 0 0, clip]{figures/part_01/p1_ch03/cosine.pdf}
    \caption{余弦函数$y=\cos(x)$在区间$[0,\pi]$上的图像。曲线自$(0,1)$单调递减，于$x=\pi/2$处过零，最终到达$(\pi,-1)$，且两端斜率为零。}
    \label{fig:cosine}
\end{figure}

这就是\textbf{余弦退火（Cosine Annealing）}的基本思想，其学习率计算公式为：
$$
\eta_t = \eta_{\min} + \frac{1}{2}(\eta_{\max} - \eta_{\min})
\left(1 + \cos\left(\frac{T_{\text{cur}}}{T_{\max}}\pi\right)\right)
$$
其中，$\eta_{\max}$和$\eta_{\min}$分别表示学习率的上限和下限，$T_{\text{cur}}$表示当前的训练轮次，$T_{\max}$表示总的训练轮次。随着训练进行，$T_{\text{cur}}$从$0$逐渐增加到$T_{\max}$，余弦函数的自变量也会从$0$逐渐变化到$\pi$，因此学习率会按照余弦曲线的变化规律，从$\eta_{\max}$连续、平滑地下降到$\eta_{\min}$。

为了更直观地理解这个过程，我们代入一组具体的参数。假设设置$\eta_{\max}=0.1$、$\eta_{\min}=0$，总训练轮次$T_{\max}=100$个Epoch，那么公式可以简化为：
$$
\eta_t = 0.05\left(1+\cos\left(\frac{T_{\text{cur}}}{100}\pi\right)\right)
$$
在训练刚开始时，$T_{\text{cur}}=0$，此时$\cos(0)=1$，因此$\eta_0=0.1$，学习率处于最大值。训练进行到四分之一，即$T_{\text{cur}}=25$时，有$\cos(\pi/4)\approx0.707$，此时学习率约为$0.085$，相比初始值仅下降了$0.015$，说明训练初期下降得比较缓慢。而当训练进行到一半，即$T_{\text{cur}}=50$时，$\cos(\pi/2)=0$，学习率下降到$0.05$。继续训练到四分之三，即$T_{\text{cur}}=75$时，$\cos(3\pi/4)\approx-0.707$，学习率约为$0.015$，此时下降速度明显加快。最后，当训练进行到$T_{\text{cur}}=100$时，$\cos(\pi)=-1$，学习率恰好下降到$\eta_{\min}=0$。

由这个例子可以看出，余弦退火不会让学习率突然下降，而是让它在训练初期下降较慢，中期较快，后期再次放缓。相比阶跃衰减，这种变化更加连续、平滑，使参数更新能够从前期的快速调整逐渐过渡到后期的精细微调。

\begin{purpleBox}{学习率预热（Warmup）}
在使用余弦退火或设置较大的初始学习率时，通常还会搭配一种常见的技术：\textbf{学习率预热（Warmup）}。在训练刚开始时，模型的参数通常是随机初始化的，梯度也可能存在较大的波动。如果一开始就直接使用较大的学习率$\eta_{\max}$，参数更新幅度可能过大，使训练初期的Loss出现明显震荡，甚至导致训练不稳定。

因此，学习率预热的核心思想，就是在训练初期的几个Epoch内，不直接使用较大的学习率，而是让学习率从一个较小的值逐渐增加到设定的最大值$\eta_{\max}$。例如，可以从$0$开始，按照线性方式逐步增加学习率。完成预热后，再进入余弦退火阶段，使学习率从$\eta_{\max}$逐渐下降到$\eta_{\min}$。

\begin{figure}[H]
    \boxedimage[width=0.6\textwidth, trim=0 0 0 0, clip]{figures/part_01/p1_ch03/warmup.pdf}
    \caption{学习率预热和余弦退火的学习率调度示意图。前$T_{\text{warm}}$个epoch内学习率从$0$线性升至$\eta_{\max}$，随后按余弦曲线平滑衰减至$\eta_{\min}$，两段在$t=T_{\text{warm}}$处连续衔接。}
    \label{fig:warmup}
\end{figure}

如\cref{fig:warmup}所示，整个训练过程可以分为两个阶段：首先通过学习率预热逐渐增大学习率，使模型平稳进入训练状态；随后通过余弦退火逐步减小学习率，使参数更新从快速调整过渡到精细微调。因此，学习率预热+余弦退火是一种常见的学习率调度组合。
\end{purpleBox}

\section{批量策略}

我们接着看梯度下降的参数更新公式：
$$
\theta \leftarrow \theta - \eta \cdot \frac{\partial \mathcal{L}}{\partial \theta}
$$
在之前的推导中，我们实际上默认了一种理想情况：每次更新参数时，模型都会使用整个训练数据集来计算平均损失，再根据得到的梯度更新参数。这样得到的梯度能够充分利用所有训练样本，因此方向最为准确。

但在实际训练中，这种做法并不总是可行。假设一个数据集包含100万张图片，如果每更新一次参数，都需要将100万张图片全部输入网络，完成一次前向传播和反向传播，那么单次参数更新的计算成本将不可接受。于是，一个问题就出现了：每次更新参数时，到底应该使用多少个样本来计算梯度？

最直接的做法，就是使用全部训练数据。如果数据集规模较小，能够一次性装入内存，那么让所有样本共同参与一次参数更新，可以得到最完整的梯度信息。这就是\textbf{全批量梯度下降（Batch Gradient Descent, BGD）}，其梯度计算公式为：
$$
\frac{\partial \mathcal{J}}{\partial \theta} = \frac{1}{n} \sum_{i=1}^{n} \frac{\partial \mathcal{L}^{(i)}}{\partial \theta}
$$
其中，$n$表示训练数据集中的样本总数。由于每次更新都使用全部样本，BGD得到的梯度估计最充分，参数更新过程也更加平稳。但与此同时，每进行一次参数更新都必须遍历整个数据集，因此数据规模越大，单次更新的计算成本就越高。

既然使用全部数据计算梯度的代价太高，那么我们也可以反过来，只使用一个样本进行参数更新。这就是\textbf{随机梯度下降（Stochastic Gradient Descent, SGD）}。每次随机选取一个样本每次随机选取一个样本$i$，直接使用该样本的梯度代表整个数据集的梯度：
$$
\frac{\partial \mathcal{J}}{\partial \theta} \approx \frac{\partial \mathcal{L}^{(i)}}{\partial \theta}
$$
这样一来，每次参数更新只需要处理一个样本，计算成本大幅降低，参数也能够更加频繁地更新。但由于单个样本只能反映整个数据集的一小部分信息，梯度估计的方差较大，参数更新会受到明显的随机波动影响。因此，SGD的优化轨迹通常比较曲折，在最优解附近也容易出现震荡。但是，这种随机性有时反而能够帮助模型跳出某些局部区域，使参数继续寻找更好的解。

全批量梯度下降计算成本太高，而随机梯度下降又容易产生较大的波动，我们可以采取一个折中方案：每次随机选取一小批样本，用这一批样本的平均梯度来表示整个数据集的梯度。这就是\textbf{小批量梯度下降（Mini-batch Gradient Descent）}，也是现代深度学习中最常用的训练方式。其梯度计算公式为：
$$
\frac{\partial \mathcal{J}}{\partial \theta} \approx \frac{1}{b} \sum_{i=1}^{b} \frac{\partial \mathcal{L}^{(i)}}{\partial \theta}
$$
其中，$b$表示每批抽取的样本数量，即Batch Size。Mini-batch的核心优势在于，它在计算效率和训练稳定性之间取得了较好的平衡。一方面，每次只使用一小部分样本计算梯度，不需要像全批量梯度下降那样遍历整个数据集，因此能够显著降低单次参数更新的计算成本；另一方面，一个Batch中包含多个样本，对它们的梯度进行平均后，又能够减小单个样本带来的随机波动，使参数更新更加稳定。因此，在现代深度学习中，通常采用Mini-batch的方式进行训练。并且，一般来说，当我们谈论“梯度下降”时，指的就是这种基于Mini-batch的梯度下降方法。

\begin{figure}[H]
    \boxedimage[width=0.9\textwidth, trim=0 0 0 0, clip]{figures/part_01/p1_ch03/comparison_of_batching_strategies.pdf}
    \caption{BGD、SGD与Mini-batch三种梯度下降轨迹在损失等高线图上的对比图。等高线由外向内递减、中心为极小值点。}
    \label{fig:comparison_of_batching_strategies}
\end{figure}

如\cref{fig:comparison_of_batching_strategies}所示，我们以损失函数的等高线图作为背景，等高线由外向内逐渐降低，中心位置为极小值点。可以看到，BGD使用全部样本计算梯度，更新方向较为稳定，因此沿着山谷平滑地向谷底逼近，轨迹近乎直线；SGD每次只使用一个样本，梯度受到单个样本的影响较大，因此更新过程中会出现明显的随机波动，轨迹也更加曲折；Mini-batch则介于两者之间，虽然仍然存在一定的抖动，但整体更新方向更加稳定，既保留了一定的随机性，又能够较快地向最优解靠近。

\begin{purpleBox}{Epoch、Batch Size和Iteration}

在深度学习的训练过程中，Epoch、Batch Size和Iteration是三个非常重要的基本概念。其中，一个\textbf{Epoch（轮次）}表示模型将整个训练数据集完整地训练一遍，也就是完成一轮训练。由于模型通常需要反复学习训练数据，才能逐渐调整出合适的参数，因此一次完整的模型训练通常需要进行多个Epoch。

\textbf{Batch Size（批大小）}表示每次参数更新时，模型同时处理的样本数量，也就是前面Mini-batch中的$b$。例如，当Batch Size设置为32时，模型每次使用32个样本计算梯度，并根据计算得到的梯度更新一次模型参数。Batch Size越小，每次更新使用的样本越少，梯度波动通常越明显；Batch Size越大，梯度估计通常更加稳定，但需要更多的计算资源和显存。

在明确了Epoch和Batch Size之后，\textbf{Iteration（迭代次数）}也就比较容易理解了。一次Iteration就是模型完成一次参数更新。因此，在一个Epoch中，模型会按照Batch Size将整个训练数据集分成若干个Batch，每处理完一个Batch，就进行一次参数更新，也就完成一次Iteration。

例如，假设训练数据集包含1000个样本，Batch Size设置为32，则完成一个Epoch需要：
$$
\left\lceil\frac{1000}{32}\right\rceil=32
$$
次Iteration。由于$32\times31=992$，因此前31个Batch各包含32个样本，最后一个Batch包含剩余的8个样本。如果总共训练50个Epoch，那么整个训练过程就需要进行：
$$
32\times50=1600
$$
次参数更新，也就是1600次Iteration。

因此，可以简单地记住：Epoch表示完整遍历数据集的次数，Batch Size表示一次使用多少个样本，Iteration表示进行多少次参数更新。
\end{purpleBox}

需要说明的是，严格来说，随机梯度下降（SGD）是指每次使用单个样本计算梯度并更新参数。但在现代深度学习中，实际训练通常以Mini-batch为单位进行参数更新。例如，\code{PyTorch}中的\code{torch.optim.SGD}虽然名称仍然使用SGD，但实际训练时也是在每个Mini-batch上计算梯度并更新参数。因此，在后面的内容中，我们提到SGD时，通常指的就是这种基于Mini-batch的随机梯度下降。

\section{动量法}

SGD（此处及后文的SGD均指Mini-batch GD，下同）能够在计算效率和训练稳定性之间取得较好的平衡。但在实际训练中，损失函数的形状往往并不规则。如\cref{fig:ill_conditioned}所示，有些损失函数会呈现出极其狭长的峡谷。

\begin{figure}[H]
    \boxedimage[width=0.95\textwidth, trim=0 0 0 0, clip]{figures/part_01/p1_ch03/ill_conditioned.pdf}
    \caption{病态曲率损失面（左）及其上的SGD参数优化轨迹（右）。左图等高线围成狭窄且弯曲的谷，垂直谷方向的曲率与梯度远大于沿谷方向；右图中SGD的更新步因此在谷壁之间来回横跳，难以沿谷方向向极小值点（白色圆圈）有效推进。}
    \label{fig:ill_conditioned}
\end{figure}

这种地形被称为\textbf{病态曲率（Ill-conditioned Curvature）}。在这种地形中，垂直于峡谷底部的方向上，损失函数变化很快，梯度往往较大；而沿着峡谷底部的方向，损失函数变化较慢，梯度相对较小。如果使用普通的Mini-batch SGD进行优化，参数更新就容易受到大梯度方向的影响，在峡谷两侧反复来回移动，而真正指向极小值的方向更新却比较缓慢。最终，参数会沿着一条不断左右摆动的“之”字形路径向前移动，虽然Loss在逐渐下降，但收敛速度却非常缓慢。

为什么会出现这种现象？一个重要原因是，普通SGD在每次更新参数时，只根据当前时刻的梯度决定更新方向，而不会保留之前的更新信息。在狭长的峡谷中，垂直于峡谷底部的方向梯度较大，因此参数很容易被推向峡谷另一侧；一旦越过谷底，梯度方向又会发生改变，于是下一次更新又会被推向另一侧。这样反复进行，参数就会在峡谷两侧不断来回摆动，而沿着峡谷底部方向的梯度较小，参数向最优解前进的速度反而很慢。那么，我们能不能让参数在更新时保留之前的运动趋势，而不是每次都只根据当前梯度决定方向？也就是让参数在考虑当前梯度的同时，也参考过去的更新方向，从而减少这种反复摆动。

那为什么考虑之前的更新方向，就能够减小这种摆动呢？这是因为在狭长峡谷中，参数在垂直于峡谷底部的方向上会不断受到方向相反的梯度影响：这一步被推向左侧，下一步又被推向右侧。如果把之前的更新方向保留下来，这些相反方向的更新会在累积过程中相互抵消，从而减弱参数在横向上的来回摆动。与此同时，沿着峡谷底部方向的梯度通常具有较强的一致性，连续多次更新会在相同方向上逐渐积累，使参数能够更快地沿着峡谷底部向最优解移动。

根据这个思路，我们可以在参数更新过程中引入一个速度变量$v_t$，用它记录历史梯度积累形成的运动趋势。这样，参数的更新就不再只取决于当前梯度，而是会同时受到之前更新方向的影响。这种利用历史梯度信息来改进参数更新的方法，就是\textbf{动量法（Momentum）}。

有了速度变量之后，我们就可以将这一思想写成数学形式：
$$
\begin{gathered}
v_t = \gamma v_{t-1} + \eta g_t \\
\theta_t = \theta_{t-1} - v_t
\end{gathered}
$$
其中，$v_{t-1}$表示上一时刻累积的速度，$\gamma$表示动量系数，用来控制历史信息保留的程度，通常取0.9；$\eta$表示学习率，$g_t$表示当前时刻的梯度，即：
$$
g_t=\nabla\mathcal{L}(\theta_{t-1})
$$
从这个公式可以看出，当前的速度$v_t$由两部分组成：一部分来自上一时刻已经积累的速度，另一部分来自当前梯度带来的更新。也就是说，参数这一步往哪里走，不再完全由当前梯度决定，而是会在历史运动趋势的基础上进行调整。

进一步展开这个递推关系，可以得到：
$$
v_t = \eta g_t + \gamma\eta g_{t-1} + \gamma^2\eta g_{t-2} + \cdots
$$
可以看到，越近的历史梯度影响越大，越早的梯度影响则会按照$\gamma$的幂逐渐减弱。因此，动量法本质上是在不断积累过去的梯度信息，同时让较久远的信息逐渐衰减。这也使得参数更新具有了一定的“惯性”。

\begin{figure}[H]
    \boxedimage[width=0.9\textwidth, trim=0 0 0 0, clip]{figures/part_01/p1_ch03/ill_conditioned_sovled.pdf}
    \caption{病态曲率损失面（左）与使用动量法的参数优化轨迹（右）。在窄而弯的谷地形上，动量法通过速度项累积历史梯度，抵消横向震荡并叠加沿谷推进；轨迹仅在初期轻微过冲，随后即沿谷底平滑收敛至极小值点（白色圆圈）。}
    \label{fig:ill_conditioned_sovled}
\end{figure}

使用动量法之后的优化结果如\cref{fig:ill_conditioned_sovled}所示。将它与普通SGD的优化轨迹进行对比可以发现，普通SGD容易在峡谷两侧反复震荡，形成明显的“之”字形路径；而加入动量后，历史梯度带来的累积效应能够减弱横向摆动，同时增强沿峡谷方向的运动，使参数更加平稳、快速地向极小值靠近。由此，动量法不仅能够缓解病态曲率带来的震荡问题，还能够利用历史信息加快收敛。

\section{自适应学习率算法}

我们通过学习率衰减让学习率随着训练过程逐渐减小，通过动量法利用历史梯度改善参数更新方向。不过，这些方法都有一个共同点：在同一个时刻，网络中的所有参数都使用相同的学习率。虽然不同参数对应的梯度大小不同，因此实际的参数更新幅度也会有所不同，但学习率本身对所有参数都是统一的。然而，在实际任务中，不同参数的梯度情况可能存在很大差异。有些参数长期处于较大的梯度区域，需要逐渐减小更新步长；而另一些参数的梯度一直较小，则可能需要更大的步长才能得到充分更新。

那么，我们能不能让不同参数根据自身的历史梯度，自动调整各自的学习率？这就是\textbf{自适应学习率算法}的基本思想。与统一设置一个全局学习率不同，它会根据每个参数过去的梯度情况，动态调整这个参数当前的有效学习率。

\subsection{AdaGrad}

最直接的做法是：既然不同参数的更新情况不同，那么就分别记录每个参数过去的梯度。如果某个参数长期收到较大的梯度，说明它已经被频繁地调整过，那么后续可以逐渐减小它的学习率；反过来，如果某个参数历史上的梯度一直较小，那么就保留较大的学习率，让它能够继续得到充分更新。那么，怎么记录一个参数过去的梯度呢？我们可以把从训练开始到当前时刻为止的梯度平方逐步累加起来，并用这个累积结果调整当前的学习率。这就是\textbf{AdaGrad（Adaptive Gradient）}的基本思想\footnote{AdaGrad由John Duchi、Elad Hazan和Yoram Singer于2011年提出，相关工作发表于论文\href{http://jmlr.org/papers/v12/duchi11a.html}{\textit{Adaptive Subgradient Methods for Online Learning and Stochastic Optimization}}。}。其更新形式可以写为：
$$
\begin{gathered}
G_t = \sum_{k=1}^{t} g_k^2 \\
\theta_t = \theta_{t-1} - \frac{\eta}{\sqrt{G_t+\epsilon}}g_t
\end{gathered}
$$
其中，$g_t$表示当前时刻的梯度，$\eta$表示基础学习率，$G_t$表示从第1步到第$t$步的梯度平方累积和，$\epsilon$是一个很小的常数（如$10^{-8}$），用于避免分母为$0$。

从公式中可以直接看出，$G_t$越大，分母就越大，当前参数的有效学习率也就越小。因此，那些历史梯度较大、更新较频繁的参数，会逐渐减小更新步长；而历史梯度较小的参数，则能够保留相对较大的学习率。

我们可以用一个简单的例子来理解这一点。假设模型中有两个参数$w_1$和$w_2$，虽然它们都会在每次参数更新时参与计算，但由于不同参数对Loss的影响不同，它们在不同训练步骤中得到的梯度大小也可能不同。假设基础学习率$\eta=1$。对于参数$w_1$，每一步的梯度都约为$1$；而对于参数$w_2$，每一步的梯度都约为$0.1$。训练进行到第$100$步时，两个参数的梯度平方累积和分别为：
$$
\begin{gathered}
G_{100}^{(1)}=100\times1^2=100 \\
G_{100}^{(2)}=100\times0.1^2=1
\end{gathered}
$$

因此，AdaGrad为两个参数计算出的有效学习率分别为：
$$
\begin{gathered}
\eta_{100}^{(1)} = \frac{1}{\sqrt{100}} = 0.1 \\
\eta_{100}^{(2)} = \frac{1}{\sqrt{1}} = 1
\end{gathered}
$$

可以看到，虽然两个参数使用的是相同的基础学习率$\eta=1$，但由于它们历史上的梯度大小不同，AdaGrad会自动为它们分配不同的有效学习率：梯度长期较大的$w_1$使用较小的学习率，而梯度较小的$w_2$则保留较大的学习率。这就是AdaGrad为不同参数自适应调整学习率的核心思想。

不过，AdaGrad也存在一个明显的问题。$G_t$是一个不断累加的量，只会增加而不会减少。随着训练进行，越来越多的历史梯度被加入其中，分母会不断变大，导致有效学习率持续下降。训练进行到后期时，即使当前梯度仍然比较明显，学习率也可能已经小到几乎无法继续更新参数，最终导致训练停滞。

\subsection{RMSProp}

面对AdaGrad的这一问题，一个比较直接的改进思路就是：没有必要永久保留所有历史梯度，只需要重点保留最近的梯度信息，让较早的梯度逐渐失去影响。

实际上，我们在前面的动量法中已经有过类似的思路：
$$
v_t = \eta g_t + \gamma\eta g_{t-1} + \gamma^2\eta g_{t-2} + \cdots
$$
可以看到，越新的梯度对当前更新的影响越大，而较早的梯度会随着时间推移逐渐减弱。现在，我们可以采用类似的方式来记录梯度平方：不再把所有历史梯度平方简单地累加起来，而是给当前梯度平方分配一定的权重，同时保留过去的信息，并让更久远的信息逐渐减弱。这样得到的量会更加关注近期的梯度变化，也不会随着训练时间的增加而无限增大。

这就是\textbf{RMSProp（Root Mean Square Propagation）}的基本思想\footnote{RMSProp并未以正式论文的形式发表，其提出与相关说明主要见Tijmen Tieleman和Geoffrey Hinton于2012年发布的课程讲义\href{https://www.cs.toronto.edu/~tijmen/csc321/slides/lecture_slides_lec6.pdf}{\textit{RMSprop: Divide the Gradient by a Running Average of its Recent Magnitude}}（见课件第26–31页；虽封面署名其他助教，但第29页注明除以梯度均方根的核心技巧源自Tijmen Tieleman，故学术界通常将其列为第一作者引用），课程名称为\textit{Neural Networks for Machine Learning}。}。它不再像AdaGrad那样将所有历史梯度平方不断累加，而是对近期的梯度平方进行加权平均：
$$
E[g^2]_t = \gamma E[g^2]_{t-1} + (1-\gamma)g_t^2
$$
其中，$g_t$表示第$t$次参数更新时计算得到的梯度，$g_t^2$表示当前梯度的平方；$E[g^2]_t$表示截至当前时刻，对近期梯度平方进行加权后得到的结果，可以理解为对参数近期梯度大小的一种记录。$E[g^2]_{t-1}$则是上一步的记录，当新的梯度平方$g_t^2$到来后，RMSProp会将旧的记录和当前的梯度平方按照不同的权重进行组合。其中，$\gamma$表示对历史信息的保留程度，通常取$0.9$。例如，当$\gamma=0.9$时，上一步的记录保留$90\%$，当前梯度平方贡献$10\%$。随着这一过程不断重复，较早的梯度会被一次又一次地乘以$\gamma$，因此其影响会逐渐减弱，而较新的梯度则能够更快地反映到当前的记录中。

我们进一步看一下$\gamma=0.9$时，这种历史信息会衰减得有多快。一个较早的梯度经过多次更新后，其影响会不断乘以$0.9$。经过$10$步后，
$$
0.9^{10}\approx0.35
$$
也就是说，它原本的影响只剩下约$35\%$；经过$20$步后，
$$
0.9^{20}\approx0.12
$$
只剩下约$12\%$；到了$30$步后，
$$
0.9^{30}\approx0.04
$$
影响已经降到约$4\%$。因此，$\gamma=0.9$会让较近的历史信息保留下来，而较久远的信息则逐渐被淡化。

得到$E[g^2]_t$之后，我们再利用它来调整当前梯度$g_t$的有效学习率：
$$
\theta_t = \theta_{t-1} - \frac{\eta}{\sqrt{E[g^2]_t+\epsilon}}g_t
$$
可以看到，$E[g^2]\_t$越大，分母就越大，当前梯度对应的有效学习率也就越小；反之，$E[g^2]\_t$越小，学习率就越大。因此，RMSProp会根据参数近期的梯度大小，自动调整它的更新步长。

这样一来，AdaGrad中“只记不忘”的累积方式就变成了“边记边忘”。由于$E[g^2]_t$采用的是近期梯度平方的加权平均，而不是像$G_t$那样不断累加，因此不会随着训练进行而无限增大，从而避免学习率不断衰减到接近于$0$，使模型在训练后期仍然能够保持有效的参数更新。不过，RMSProp解决的主要是每个参数应该走多远的问题，而没有解决应该往哪个方向走的问题。它通过历史梯度平方来调整参数的更新步长，但更新方向仍然主要由当前梯度$g_t$决定。因此，如果当前Batch带来的梯度存在较大的随机波时，参数的更新方向也可能随之发生变化，仍然容易受到当前梯度噪声的影响。

\subsection{Adam}

既然动量法利用历史梯度的信息，解决了参数往哪个方向走的问题；RMSProp则根据梯度大小调整参数的更新步长，解决了参数走多远的问题。那我们能不能把这两种方法结合起来？既利用历史梯度来平滑更新方向，又根据每个参数自身的梯度情况自动调整学习率。这就是\textbf{Adam（Adaptive Moment Estimation）}的基本思想\footnote{Adam由Diederik Kingma和Jimmy Ba于2015年提出，相关工作发表于论文\href{https://arxiv.org/abs/1412.6980}{\textit{Adam: A Method for Stochastic Optimization}}。}。

Adam同时维护两个量：一阶矩$m_t$和二阶矩$v_t$\footnote{一阶矩和二阶矩来自统计学中的\textbf{矩（Moment）}：一阶原点矩为$E[X]$，二阶原点矩为$E[X^2]$。在Adam中，$m_t$是对梯度一阶矩的指数加权估计，用于反映梯度的平均方向；$v_t$是对梯度二阶矩的指数加权估计，用于反映梯度的尺度。需要注意的是，$v_t$对应的是$E[g_t^2]$，而不是方差。}。一阶矩记录梯度的加权平均，用于平滑更新方向；二阶矩记录梯度平方的加权平均，用于调整参数的有效学习率。其核心公式为：
$$
\begin{gathered}
m_t = \beta_1 m_{t-1} + (1-\beta_1)g_t \\
v_t = \beta_2 v_{t-1} + (1-\beta_2)g_t^2
\end{gathered}
$$
其中，$\beta_1$控制历史梯度信息的保留程度，$\beta_2$控制历史梯度平方信息的保留程度。可以看到，第一个公式与动量法的思想类似，通过保留历史梯度来平滑当前的更新方向；第二个公式与RMSProp类似，通过记录近期梯度平方的大小来调整更新步长。

不过，$m_t$和$v_t$在训练开始时通常都初始化为$0$。这会带来一个问题：在最初几次更新中，它们积累的梯度信息还比较少，因此得到的数值会明显偏小，不能很好地反映当前梯度的真实情况。例如，假设只有第$1$次更新时产生了梯度，之后暂时不考虑新的梯度，那么由于$m_0=0$，第一次计算得到的$m_1$只保留了当前梯度的一部分，而不是完整的梯度信息。随着训练继续进行，这种从$0$开始累积带来的低估会逐渐减弱。为了解决这个问题，Adam引入了\textbf{偏差修正（Bias Correction）}：
$$
\hat{m}_t = \frac{m_t}{1-\beta_1^t}, \qquad \hat{v}_t = \frac{v_t}{1-\beta_2^t}
$$
这里的$1-\beta_1^t$和$1-\beta_2^t$就是用来补偿初始化为$0$所造成的低估。训练刚开始时，$\beta_1^t$和$\beta_2^t$还比较接近$1$，分母较小，因此修正幅度较大；随着$t$不断增加，$\beta_1^t$和$\beta_2^t$逐渐接近$0$，修正作用也就越来越小。

得到修正后的$\hat m_t$和$\hat v_t$后，Adam最终的参数更新公式为：
$$
\theta_t = \theta_{t-1} - \frac{\eta}{\sqrt{\hat{v}_t}+\epsilon}\hat{m}_t
$$

整个过程可以理解为：首先根据当前梯度$g_t$更新一阶矩和二阶矩；然后对这两个统计量进行偏差修正；最后利用修正后的一阶矩确定参数更新方向，再利用二阶矩调整更新步长。于是，Adam就把“方向”和“步长”两个问题统一到了同一个更新公式中。

为了更直观地理解Adam的作用，我们还是回到前面的峡谷问题。假设某个参数在优化过程中不断来回摆动，得到的梯度依次为：
$$
g_1=1,\qquad g_2=-1,\qquad g_3=1,\qquad g_4=-1
$$
可以看到，梯度方向不断发生变化。此时，一阶矩会对这些历史梯度进行加权平均，正负交替的梯度会在累积过程中相互抵消，因此一阶矩的变化会比原始梯度平缓得多，从而减小单次梯度波动对更新方向的影响。

而与此同时，这些梯度的绝对值始终为$1$，说明这个参数一直处于梯度较大的区域。因此，二阶矩会持续记录较大的梯度平方，并通过分母$\sqrt{\hat{v}_t}$压缩有效学习率，从而避免参数在梯度较大的区域进行过大的更新。

所以，Adam实际上同时解决了两个问题：一阶矩负责往哪个方向走，二阶矩负责走多远。一阶矩继承了动量法的思想，二阶矩则继承了RMSProp的思想，而偏差修正则进一步改善了训练初期的统计估计。

从整个优化算法的发展过程来看，可以看到一条清晰的演进路线：动量法利用历史梯度平滑更新方向，减少参数震荡；AdaGrad进一步让不同参数拥有不同的有效学习率，但由于历史梯度平方不断累积，学习率可能逐渐衰减过小；RMSProp让久远的历史信息逐渐衰减，解决了AdaGrad学习率持续下降的问题；最终，Adam将动量法和RMSProp的思想结合起来，同时对更新方向和步长进行调整。不过，Adam并非在所有任务中都一定优于其他优化器，因此实际训练时仍需要根据具体任务选择合适的方法。

\begin{quote}
    我们已经掌握了多种优化算法，能够让模型更加高效、稳定地更新参数。那么，在实际训练深度神经网络时，还会遇到其他问题吗？
\end{quote}
